% ============================================================
%  第 6 章  总结与展望
% ============================================================

\chapter{总结与展望}

% ================================================================
\section{工作总结}

本文围绕"面向多平台社交媒体的实时舆情分析与热度趋势监测"这一核心问题，完整设计并实现了一套端到端的准实时舆情监测系统。主要工作成果概括如下：

\begin{enumerate}
  \item \textbf{多平台爬虫工程化实现。}基于 MediaCrawler 框架，采用 Chrome DevTools Protocol（CDP）接管机制实现了对微博、B 站、抖音、小红书、知乎、快手、贴吧七个主流社交媒体平台的关键词评论采集，支持可配置轮询间隔的持续采集模式，并以 JSONL 格式输出，与 Kafka 消息管道实现低耦合对接。

  \item \textbf{Spark Structured Streaming 消费管道设计。}构建了基于 60 秒微批次触发的 Kafka-Spark 消费管道，实现了完整的 JSON Schema 解析、关键字段非空校验与容错处理流程，通过 Checkpoint 机制保障了异常重启后的数据不丢失。

  \item \textbf{StructBERT 情感分析服务的流处理集成。}提出了懒初始化单例模式 + 批量推理 + 异常降级的三层工程化方案，将 StructBERT 中文情感分类模型嵌入 \code{foreachBatch} 回调，实现了批量吞吐量约 32 条/秒（\code{batch\_size=32}）的高效推理，相较逐条推理提升约 5 倍。

  \item \textbf{多维度热度评分模型。}设计并实现了综合点赞数、转发数与情感极性加成的热度评分公式，以及基于热度分数与情感标签组合的三级风险评估机制，实验验证了各风险等级热度分布的单调分离性。

  \item \textbf{Delta Lake Medallion 三层湖仓架构。}在单机环境下完整实现了 Bronze（追加）—Silver（Merge 去重）—Gold（全量重算）三层 Delta Lake 数据湖架构，72 小时内处理 18,437 条评论，总存储开销约 46~MB，各端点平均响应时间低于 1.5~秒。

  \item \textbf{FastAPI 数据服务与 Vue~3 可视化大屏。}设计了 9 个 RESTful 端点，实现了基于 \code{deltalake} Python 库的无 JVM 直读方案；前端基于 Vue~3 Composition API 与 ECharts~6 实现了支持多话题切换、双粒度趋势图、评论明细与告警面板的交互式可视化大屏，并通过 CSS 媒体查询完成了多分辨率自适应适配。
\end{enumerate}

% ================================================================
\section{不足之处}

尽管本系统在功能完整性与端到端可用性方面达到了预期设计目标，但仍存在以下值得进一步改进的局限性：

\begin{enumerate}
  \item \textbf{情感模型的领域适应性不足。}当前系统使用的 StructBERT 通用情感分类模型在处理网络流行语、混合语言表达、反讽语气等非标准社交媒体文本时，准确率下降明显。由于缺乏高质量的领域标注语料，本系统未进行针对性微调，这是情感分析模块最核心的改进空间。

  \item \textbf{单机部署的可扩展性瓶颈。}\code{local[*]} 部署模式下，Spark Driver 同时承载 NLP 推理与数据写入两类计算任务，资源竞争较为突出。当监测话题数量显著增加或批次数据量急剧扩大时，单机资源将成为明显瓶颈。

  \item \textbf{API 全量数据加载问题。}\code{/api/comments} 接口当前将 Silver 表全量加载为 Pandas DataFrame 后在内存中进行过滤，当数据规模超过百万条时将出现显著的性能退化。

  \item \textbf{爬虫长期稳定性约束。}各平台反爬策略持续迭代升级，CDP 接管模式在高频采集场景下存在会话失效的风险，系统的长期无人值守运行稳定性有待进一步验证与加固。
\end{enumerate}

% ================================================================
\section{未来展望}

针对上述不足，未来工作可从以下几个方向展开：

\paragraph{情感分析能力升级}引入 Llama-3 或 Qwen 等大语言模型（LLM）进行零样本（Zero-Shot）情感分析，或构建社交媒体领域标注数据集对 StructBERT 进行指令微调，以提升对网络流行语与隐性负面情绪的识别能力。

\paragraph{分布式流处理扩展}将部署模式升级为 Spark Standalone 集群或接入云厂商托管的 EMR 服务，实现 Executor 的弹性扩缩容；同时将 StructBERT 推理服务独立化为 gRPC 微服务，彻底解除推理计算与 Spark 调度之间的资源耦合。

\paragraph{实时存储层优化}引入 Apache Hudi 的 Merge-On-Read 增量写入机制，替代当前 Gold 层全量 Overwrite 方案，降低批次写入放大系数；对 \code{/api/comments} 接口改造为基于 DuckDB 的 Parquet 直查模式，规避全量加载的内存瓶颈。

\paragraph{舆情预测与知识图谱}在现有情感分析基础上，引入时序预测模型（如 LSTM 或 Transformer 时序模块）对热度趋势进行短期预测；结合命名实体识别（NER）构建舆情知识图谱，实现话题间关联分析。

\paragraph{多模态内容分析}当前系统仅分析文本评论，未来可扩展至视频弹幕、表情包文本乃至图片内容的多模态融合分析，进一步提升舆情感知的全面性。

% ================================================================
% 致谢（不编入章节编号，加入目录）
\chapter*{致\quad 谢}
\addcontentsline{toc}{chapter}{致谢}
\markboth{致谢}{致谢}

本文的完成离不开众多人士的关怀与支持，谨在此致以诚挚的谢意。

首先，衷心感谢指导教师\textbf{陈志}老师在本课题研究过程中给予的悉心指导与宝贵建议。陈老师以严谨的治学态度、深厚的学术积淀和对工程实践的深刻洞察，在研究方向确定、技术方案论证与论文写作规范等各个环节给予了方向性的引领与细致入微的修改意见，是本文得以顺利完成的最重要支撑。

其次，感谢天津工业大学计算机科学与技术学院为本人提供了良好的学习与研究环境，感谢大数据 2201 班同学在日常学习与项目讨论中给予的帮助与启发。

再次，本文的实现参考和使用了 Apache Spark、Delta Lake、Apache Kafka、ModelScope StructBERT、MediaCrawler、FastAPI、Vue.js、ECharts 等优秀开源项目，谨向这些项目的所有贡献者表示敬意与感谢，开源精神深刻影响了笔者对软件工程的认知。

最后，最深切的感谢献给家人多年来无怨无悔的默默支持与鼓励，是你们的理解与陪伴赋予了坚持走到今天的力量。

由于本人知识储备与实践经验尚存不足，论文中难免存在疏漏与不当之处，诚恳地请各位专家、评委和读者批评指正。

\vspace{2em}
\begin{flushright}
  高楠\\
  2026 年 6 月于天津
\end{flushright}
